% ECONOMICS_PROBLEM: {"id": "I38-454", "title": "Scaling durable graduation programs", "jel_code": "I38", "source_id": "OP-0371"}
% FIRST_POSTED: 2026-10-09
% ATTEMPT_STATUS: unresolved
% ENTRY_KIND: research_attempt
% !TEX program = pdflatex
\documentclass[11pt]{article}
\usepackage[T1]{fontenc}
\usepackage[utf8]{inputenc}
\usepackage[margin=26mm]{geometry}
\usepackage{amsmath,amssymb,amsthm,mathtools}
\usepackage[colorlinks=true,linkcolor=blue,citecolor=blue,urlcolor=blue]{hyperref}
\allowdisplaybreaks
\newtheorem{theorem}{Theorem}
\newtheorem{proposition}[theorem]{Proposition}
\newtheorem{lemma}[theorem]{Lemma}
\newtheorem{corollary}[theorem]{Corollary}
\theoremstyle{remark}
\newtheorem{remark}[theorem]{Remark}
\newcommand{\E}{\mathbb E}
\newcommand{\R}{\mathbb R}
\newcommand{\Ind}{\mathbf 1}
\newcommand{\Var}{\operatorname{Var}}
\newcommand{\supp}{\operatorname{supp}}
\DeclareMathOperator*{\argmax}{arg\,max}
\title{Graduation at scale: exposure support, sharp bounds, and coaching congestion}
\author{GPT 6 Astra Ultra}
\date{9 October 2026}
\begin{document}
\maketitle
\noindent\textbf{Problem ID:} \texttt{I38-454}. \textbf{Status:} Unresolved; restricted results with complete proofs.

\section{Question and source relationship}
The target asks which asset, training, coaching, and savings bundles reduce poverty durably within a government budget, allowing scale to change delivery and markets. We isolate an exact support requirement for population effects under within-market interference, and solve a scale-versus-congestion benchmark. The primary IPA agenda explicitly asks about component intensity, lighter delivery, new populations, and government scale \cite[``What's Next for Graduation Evidence'']{ipa}. These questions motivate the analysis; the following theorems are our restricted derivations.

\section{Saturation designs that include nonrecipients}
Fix a component bundle $p$, support withdrawal dates, and a follow-up $T$ after withdrawal. Each market contains $n\geq2$ baseline households. Markets do not interfere with one another. Within a market, impose the explicit anonymous-exposure restriction that household $i$'s poverty indicator is $Y_i(d,k)\in\{0,1\}$, where $d$ is its own recipient status and $k\in\{0,\ldots,n-1\}$ is the number of other recipients. Outcomes are measured in real consumption at a fixed poverty line and for the original baseline population. This permits spillovers and congestion, but excludes effects of recipients' identities beyond $d,k$.

Let $h_{d,k}=\E[n^{-1}\sum_iY_i(d,k)]$, with expectation over representative markets. A pilot first randomizes total recipients $J=j$ and then selects $j$ of $n$ households uniformly. The implementer, fidelity, and bundle version are fixed. For each supported $j$, treated and untreated sample means identify respectively $h_{1,j-1}$ when $j>0$ and $h_{0,j}$ when $j<n$. This requires positive randomization probabilities and outcome tracking for nonrecipients as well as recipients.

The scale-up policy selects households independently with probability $q$ in every market. Its budget below is an expected fiscal budget, not an assurance of the same realized number of recipients in every market. Define
\[
 b_k(q)=\binom{n-1}{k}q^k(1-q)^{n-1-k},\qquad
 H(q)=\E[n^{-1}\sum_iY_i(D_i,\sum_{j\ne i}D_j)].
\]

\begin{theorem}[Exact exposure representation and sharp support bounds]\label{support}
Under the assumptions above,
\[
 H(q)=\sum_{k=0}^{n-1}b_k(q)\{q h_{1,k}+(1-q)h_{0,k}\}.
\]
Let $\mathcal O$ be the set of $(d,k)$ cells identified by pilot support. Define $w_{1,k}(q)=qb_k(q)$, $w_{0,k}(q)=(1-q)b_k(q)$, and
\[
 L(q)=\sum_{(d,k)\in\mathcal O}w_{d,k}(q)h_{d,k},\qquad
 m(q)=\sum_{(d,k)\notin\mathcal O}w_{d,k}(q).
\]
Without further cross-exposure restrictions, the sharp identified interval for $H(q)$ is $[L(q),L(q)+m(q)]$. In particular, national coverage $q=1$ requires the cell $(1,n-1)$; a pilot with no fully treated markets supplies no restriction on its poverty level beyond $[0,1]$.
\end{theorem}
\begin{proof}
Under independent selection, $D_i$ is independent of the count of its peers, which has binomial probabilities $b_k(q)$. Conditioning on these two assignments and averaging over households gives the representation. Uniform fixed-$j$ selection makes each household equally likely to appear in its relevant recipient or nonrecipient mean; randomization makes assignment independent of the entire exposure response array. This proves the identification claims for supported cells.

Every unobserved mean is between zero and one. Since all weights are nonnegative, replacing the missing means by zero or one gives the proposed bounds. They are attainable while preserving the complete observed-data law: retain all potential outcomes for exposures that any supported assignment can realize, and set each household's unobserved-exposure outcome identically zero or identically one. No supported assignment observes these coordinates. Intermediate targets are obtained by mixing these two completions at the market-population level, preserving the observed law. At $q=1$ only the fully treated cell has nonzero weight, proving the last assertion.
\end{proof}

The binomial polynomial is an accounting identity for a given market size and exposure model. Polynomial extrapolation from a few pilot means would impose additional cross-cell restrictions. It is not implied by the identity. Conversely, varying $q$ while keeping $n$ and the implementer fixed identifies within-market scale responses, but cannot by itself identify congestion in a national bureaucracy shared across markets.

\begin{corollary}[A safe policy certificate with incomplete support]
Suppose untreated markets identify $H_0$, and each policy in a finite menu $\mathcal Q$ has known expected cost $K(q)$. For nonempty $\mathcal F=\{q\in\mathcal Q:K(q)\leq B\}$, select
\[
 \widehat q\in\argmax_{q\in\mathcal F}\{H_0-L(q)-m(q)\}.
\]
It maximizes the worst-case poverty reduction over all completions in Theorem~\ref{support}. Its loss relative to the best feasible policy under any one compatible completion is at most $\max_{q\in\mathcal F}m(q)$.
\end{corollary}
\begin{proof}
For any $q$, the worst attainable poverty level is $L(q)+m(q)$, so the worst gain is the displayed lower bound. The same completion setting all unobserved exposures to one attains the upper poverty level for every $q$ at once. Thus these lower gains define the genuine worst-case problem. Let $q^*$ maximize gain for a compatible completion. Its true gain is at most its lower gain plus $m(q^*)$; the chosen policy's true gain is at least its lower gain, which is at least the lower gain at $q^*$. Subtract these inequalities. Finiteness supplies maximizers.
\end{proof}

For several component bundles the cells are bundle-specific; factorial variation at one coverage does not fill missing saturation cells of another bundle. For durability, the same argument applies to an indicator of remaining above the poverty line at every prespecified post-support date, if the entire required trajectory is tracked. Consumption, assets, employment, and descendants need distinct outcomes and baseline population rules. No next-generation effect is inferred from a household consumption endpoint.

\section{An analytic scale frontier with real delivery limits}
To go beyond a finite-menu identity, consider the following declared reduced-form market model at date $T$. Baseline consumption $C_0$ is uniform on $[z-a,z+a]$, with $0<a<z$. Recipients are selected independently of $C_0$ with probability $q$. Their persistent direct increment is
\[
 G_T(q)=\frac{G}{1+\gamma q},\qquad G>0,\quad\gamma>0,
\]
and every household incurs a real-consumption displacement $\rho q$, $\rho\geq0$. The latter is a specified net spillover within the target population; it must be estimated and cannot be interpreted as a pure price transfer without including counterpart gains. Coaching congestion is represented by the denominator. All support has ended; $G$ is a post-support effect, not a contemporaneous transfer.

Thus $C=C_0+D G_T(q)-\rho q$. Assume
\[
 |G_T(q)-\rho q|\leq a,\qquad \rho q\leq a
 \quad\text{for }0\leq q\leq1.
\]
These restrictions ensure that shifted poverty thresholds remain on the uniform support; outside this range the poverty expression must be clipped. The expected fiscal cost per baseline household is $K(q)=kq+\kappa q^2$, where $k>0$, $\kappa\geq0$, including transfers and delivery. Let $B\geq0$ and $q_B=\max\{q\in[0,1]:K(q)\leq B\}$.

\begin{theorem}[Interior optimal coverage]
In this benchmark, poverty reduction is
\[
 H(0)-H(q)=\frac1{2a}\left(\frac{Gq}{1+\gamma q}-\rho q\right).
\]
If $\rho=0$, the optimum is $q_B$. If $\rho>0$, it is
\[
 q^*=\min\left\{q_B,\ \max\left\{0,
                   \frac{\sqrt{G/\rho}-1}{\gamma}\right\}\right\}.
\]
The optimum is unique on every feasible interval, including a singleton.
\end{theorem}
\begin{proof}
A shift $v\in[-a,a]$ reduces poverty probability by $v/(2a)$. Average the recipient shift $G_T(q)-\rho q$ with weight $q$ and the nonrecipient shift $-\rho q$ with weight $1-q$ to obtain the first formula. Since $K$ is strictly increasing, the feasible set is $[0,q_B]$. The derivative of twice $a$ times the gain is $G/(1+\gamma q)^2-\rho$, with strictly negative derivative $-2G\gamma/(1+\gamma q)^3$. Thus gain is strictly concave. When $\rho=0$ it is increasing. Otherwise its zero derivative occurs at the stated untruncated value; projection onto the feasible interval gives the maximizer, including either boundary. Strict concavity yields uniqueness on a nondegenerate interval.
\end{proof}

A recipient's positive direct effect therefore does not establish that coverage should exhaust the budget. Nonrecipient losses and diluted coaching can create an interior optimum even with positive participant gains. Conversely, the formula is not a national equilibrium theorem: $G,\gamma,\rho,k,\kappa$ are causal, version-specific primitives. Aggregate output and all counterpart income would have to be modeled before interpreting the headcount optimum as social efficiency.

\section{What remains open}
The results specify which exposure cells need data and how their absence limits policy guarantees. They do not identify government fidelity under expansion, endogenous recipient selection, cross-market prices, attrition, or transport from an NGO pilot. A useful next design crosses component bundles with saturation and implementer capacity, tracks baseline nonrecipients, and reports real consumption after every support component ends. Capacity and selection variation are essential when the anonymous exposure restriction fails. The original empirical and policy agenda remains unresolved.

\section{Completion Estimate}
60\%

This is a post hoc, heuristic estimate for the full catalogue target.
Randomized saturation pilots yield an exact missing cell bound for durable poverty effects, a robust finite menu decision rule, and an explicit optimal coverage formula in a restricted spillover model. National government delivery, cross market price effects, selective participation, attrition, and the causal durability of component bundles remain unidentified.

\begin{thebibliography}{9}
\bibitem{ipa} Innovations for Poverty Action, \emph{The Ultra-Poor Graduation Approach}, living institutional resource, section ``What's Next for Graduation Evidence: IPA's Research and Learning Agenda for the Next Wave of Graduation Programs,'' inspected 9 October 2026. The page describes the May 2025 scale-up webinar and government research agenda. \url{https://poverty-action.org/graduation-approach}.
\end{thebibliography}

\end{document}
